% !TeX root = ../main.tex

% ==================================================
% 1.6. URL vs Local Path
% ==================================================

\subsection{URL vs Local Path}

\subsubsection{먼저 둘의 차이부터 이해하자}

지금까지 우리가 다루었던 경로는 대부분 이런 형태였다.

\begin{dhkcode}
from pathlib import Path

path = Path("data/train.csv")
\end{dhkcode}

또는:

\begin{dhkcode}
path = Path("/Users/user/project/data/train.csv")
\end{dhkcode}

이것들은 \zsh{Local Path(로컬 파일 경로)}이다.

반면 다음은 URL 이다:

\begin{dhkcode}
https://example.com/data/train.csv
\end{dhkcode}

둘 다 파일의 위치를 나타내는 것처럼 보인다.

하지만 Python에서 접근하는 방식은 완전히 다르다.

\begin{dhktable}
|구분|Local Path|URL|
|:---|:---|:---|
|목적|파일 시스템의 위치 지정|리소스의 위치 및 접근 방법 지정|
|예시|data/train.csv|https://example.com/ data/train.csv|
|접근 방식|파일 시스템|네트워크(HTTP 등)|
|대표 도구|pathlib, open()|requests, urllib|
|접근 대상|로컬 또는 마운트된 파일 시스템|웹 서버 등의 리소스|
\end{dhktable}

여기서 중요한 것은 \textbf{URL은 단순히 인터넷에 있는 파일의 경로만을 의미하지 않는다는 것}

URL은 \zsh{Uniform Resource Locator}의 약자로,

리소스에 접근하기 위한 위치와 방법을 나타낸다.

\subsubsection{Local Path는 파일 시스템에서 해석된다}

\begin{dhkcode}
from pathlib import Path

path = Path("data/train.csv")

with path.open("r", encoding="utf-8") as f:
    data = f.read()
\end{dhkcode}

\begin{dhksmry}[centered]
Python
\da
Path("data/train.csv")
\da
현재 작업 디렉터리(CWD) 기준으로 경로 해석
\da
운영체제의 파일 시스템
\da
파일 열기
\da
file object
\da
read()
\da
Python str
\end{dhksmry}

예를 들어 CWD가:

\begin{dhkzsh}
/Users/user/my_project
\end{dhkzsh}

라면:

\begin{dhkcode}
Path("data/train.csv")
\end{dhkcode}

는 다음 파일을 가리키는 상대 경로로 해석된다.

\begin{dhkzsh}
/Users/user/my_project/data/train.csv
\end{dhkzsh}

여기서 중요한 것은 \textbf{open()이 파일 시스템을 대상으로 작동한다는 것}이다.

\subsubsection{URL은 네트워크 요청을 통해 접근한다}

이번에는:

\begin{dhkzsh}
https://example.com/data/train.csv
\end{dhkzsh}

를 보자.

이 URL을 분해하면 다음과 같다.

\begin{dhkdir}
https:/example.com/data/train.csv
> https -> scheme
> example.com -> host
> /data/train.csv -> path
\end{dhkdir}

각각의 의미는 다음과 같다.

\begin{dhksmry}
    1. Scheme
        https
        어떤 방식으로 리소스에 접근할지 나타낸다.
        여기서 HTTPS 프로토콜을 사용한다는 뜻이다.


    2. Host
        example.com
        어느 서버에 접근할지 나타내는 호스트 이름이다.


    3. Path
        /data/train.csv
        서버에 요청할 리소스의 경로이다.
        여기서 주의할 점은,
        URL의 Path는 로컬 파일시스템의 path와 동일한 개념이 아니다.
    
    서버가 이 경로를 실제 파일에 연결할 수도 있지만,
    데이터베이스 조회나 프로그램 실행 결과에 연결할 수도 있다.
\end{dhksmry}

예를 들어:

\begin{dhkzsh}
https://api.example.com/users/123
\end{dhkzsh}

은 서버의 \zsh{users/123}이라는 실제 파일을 의미하는 것이 아닐 수 있다.

서버가 사용자 123의 정보를 조회해서 JSON을 반환하는 API일 수도 있다.

\subsubsection{그래서 \zsh{open()}으로 URL을 직접 열 수 없다}

다음 코드를 보자.

\begin{dhkcode}
with open("https://example.com/data/train.csv") as f:
    data = f.read()
\end{dhkcode}

이 코드는 일반적으로 실패한다.

왜일까?

\zsh{open()}은 전달된 문자열을 \textbf{파일 시스템 경로}로 해석하기 때문이다.

즉 Python은 이것을 네트워크 주소라고 판단해서 HTTP 요청을 보내지 않는다.

반면 HTTP 요청을 보내려면 다음과 같은 라이브러리를 사용할 수 있다.

\begin{dhkcode}
import requests

url = "https://example.com/data/train.csv"

response = requests.get(url, timeout=10)
response.raise_for_status()

data = response.text
\end{dhkcode}

전체 흐름은 다음과 같다.

\begin{dhksmry}[centered]
Python
\da
requests.get(url)
\da
HTTP 요청
\da
네트워크
\da
웹 서버
\da
HTTP 응답
\da
response
\da
response.text
\da
Python str
\end{dhksmry}

여기서 \zsh{requests.get()}은 HTTP GET 요청을 보내는 함수다.

\zsh{response}는 서버에서 돌아온 HTTP 응답을 표현하는 객체이다.

\zsh{response.text}는 응답 본문을 문자로 디코딩해서 제공한다.

즉, encoding/decoding 개념과 연결된다.

\subsubsection{그런데 Pandas에서는 URL을 바로 사용할 수 있다}

우리가 지금까지 배운:

\begin{dhkcode}
pd.read_csv('data/train.csv')
\end{dhkcode}

는 로컬 파일을 읽는 코드이다.

그런데 Pandas는 다음과 같은 코드도 지원한다.

\begin{dhkcode}
import pandas as pd

url = "https://example.com/data/train.csv"

df = pd.read_csv(url)
\end{dhkcode}

단, 실제로 해당 URL에서 CSV 데이터를 제공해야 정상적으로 작동한다.

그렇다면 왜 Pandas에서는 URL을 직접 사용할 수 있을까?

\textbf{Pandas가 URL을 인식하고 내부에서 필요한 데이터 접근 및 파싱 작업을 처리하기 때문이다}.

개념적으로:

\begin{dhksmry}[centered]
pd.read_csv("data/train.csv")
\da
Local Path
\da
파일 시스템
\da
CSV
\da
DataFrame
\end{dhksmry}

반면:

\begin{dhksmry}[centered]
pd.read_csv("https://...")
\da
URL
\da
HTTP 요청
\da
CSV 데이터
\da
DataFrame
\end{dhksmry}

즉 사용자에게는 동일한 \zsh{read_csv()}인터페이스를 제공하지만,

내부적으로 데이터를 가져오는 방법은 달라질 수 있다.

이것을 \textbf{추상화(Abstraction)}의 사례로 볼 수 있다.

\subsubsection{Hugging Face에서도 비슷한 일이 일어난다}

\begin{dhkcode}
from transformers import AutoTokenizer

tokenizer = AutoTokenizer.from_pretrained(
    "distilbert-base-uncased"
)
\end{dhkcode}

여기서:

\begin{dhkzsh}
"distilbert-base-uncased"
\end{dhkzsh}

는 일반적인 URL도 아니고, 파일 경로라고 명시한 것도 아니다.

이것은 Hugging Face 모델을 식별하는 이름으로 사용할 수 있다.

라이브러리는 이 식별자를 바탕으로 필요한 파일을 찾고,

로컬 캐시에 없으면 원격 저장소에서 다운로드할 수 있다.

개념적으로:

\begin{dhksmry}[centered]
AutoTokenizer.from_pretrained(...)
\da
모델 식별자 해석
\da
로컬 캐시 확인
\da
필요한 파일이 없다면
\da
원격 저장소 접근
\da
필요한 파일 다운로드
\da
로컬 캐시
\da
Tokenizer 생성
\end{dhksmry}

물론 로컬 디렉터리를 직접 지정하는 것도 가능하다.

\begin{dhkcode}
AutoTokenizer.from_pretrained(
    "./my_tokenizer"
)
\end{dhkcode}

이 경우에는 해당 로컬 디렉터리에서 필요한 파일을 찾도록 사용할 수 있다.

따라서 \zsh{from_pretrained()}같은 함수는 단순한 파일 읽기 함수보다

더 높은 수준의 인터페이스를 제공하는 것이다.

\subsubsection{URL의 데이터도 결국 메모리나 파일로 가져올 수 있다}

예를 들어 웹에서 JSON을 가져온다고 하자.

\begin{dhkcode}
import requests

url = "https://api.example.com/config"
response = requests.get(url, timeout=10)
response.raise_for_status()
config = response.json()
\end{dhkcode}

여기서:

\begin{dhkcode}
response.json()
\end{dhkcode}

은 HTTP 응답 본문의 JSON 데이터를 Python 객체로 변환한다.

앞에서 배운:

\begin{dhkcode}
json.loads(text)
\end{dhkcode}

과 비슷한 역할을 한다고 이해하면 된다.

다만 \zsh{response.json()}은 \zsh{requests}가 제공하는 메서드다.

전체 흐름은:

\begin{dhksmry}[centered]
URL
\da
HTTP GET
\da
HTTP response
\da
JSON 데이터
\da
response.json()
\da
Python dict
\end{dhksmry}

이다.

JSON 최상위 구조에 따라 list 등 다른 Python 객체가 반환될 수 있다.

여기서 중요한 점은 \textbf{네트워크에서 데이터를 가져온다고 해서 반드시 로컬 파일로 저장해야 하는 것은 아니라는 것}이다.

응답 데이터를 바로 메모리에서 처리할 수도 있다.

\subsubsection{지금까지 배운 개념을 통합해보자}

우리가 지금까지 학습한 흐름을 연결하면 다음과 같다.

\begin{dhktree}
Python Program
>Local Path
>>pathlib
>>>File System
>>>>file object
>>>>>데이터 해석
>>>>>>Text :: str
>>>>>>JSON :: Dict/List
>>>>>>CSV :: DataFrame
>URL
>>requests
>>>HTTP 요청
>>>>HTTP response
>>>>>데이터 해석
>>>>>>Text :: str
>>>>>>JSON :: Dict/List
>>>>>>CSV :: DataFrame
\end{dhktree}

이 그림에서 한 가지를 이해하면 된다.

\textbf{데이터가 어디에 있는지와 데이터를 어떤 형식으로 해석하는지는 서로 다른 문제다}.

예를 들어 CSV는 로컬 파일에서도 읽을 수 있고 HTTP 서버에서도 가져올 수 있다.

하지만 어느 쪽에서 가져왔든 CSV 데이터를 DataFrame으로 변환하는 작업 자체는 별개의 단계이다.

이 구분은 앞으로 데이터 파이프라인을 설계할 때도 중요하다.

\subsubsection{핵심 정리}

기억할 것은 4 가지이다.

\begin{dhksmry}
    1. Local Path는 파일 시스템에서 해석되는 경로이며,
        pathlib 와 open() 으로 다룬다.
    

    2. URL은 리소스의 위치와 접근 방법을 나타내며,
        HTTP URL은 requests 같은 도구로 접근할 수 있다.
    

    3. pd.read_csv()처럼 고수준 라이브러리는 
        로컬 파일과 URL을 모두 지원할 수 있다.


    4. 데이터에 접근하는 단계와 JSON·CSV 등의 형식을 해석하는 단계는 구분된다.
\end{dhksmry}

\subsubsection{HTTP 요청과 응답}

HTTP는 클라이언트와 서버가 요청(request)과 응답(response)을 주고받기 위한 통신 규칙이다.

예를 들어 Python에서 공공데이터를 요청한다고 생각해보자.

\begin{dhkbranch}[in=8]
node main :: Python 프로그램 (Client)
arrow main :: HTTP Request
node main :: 공공데이터 서버
arrow main :: HTTP Response
node main :: Python 프로그램
\end{dhkbranch}

이때 Python 프로그램이 클라이언트다.

공공데이터를 제공하는 쪽이 서버이다.

핵심은 \textbf{HTTP 통신에서는 클라이언트가 요청을 보내고 서버가 응답한다}는 구조다.

\subsubsection{HTTP Request는 어떻게 구성될까?}

HTTP 요청을 개념적으로 분해하면 다음과 같다.

\begin{dhkdir}
HTTP Request
>Method
>>GET / POST / PUT / DELETE ...
>URL
>>https://api.example.com/data
>Headers
>>요청에 대한 부가 정보
>Body
>>서버에 전달할 데이터
\end{dhkdir}

하나씩 보자.

\subsubsection{Method - 무엇을 하려는가?}

HTTP Method는  서버에 요청하는 작업의 성격을 나타낸다.

대표적인 두 가지가 있다.

\textbf{GET}

서버에서 데이터를 가져오는 데 사용한다.

\begin{dhkcode}
response = request.get(
    "https://api.example.com/data"
)
\end{dhkcode}

예를 들어:

\begin{dhksmry}
    - 공공데이터 조회

    - 특정 금융상품 정보 조회

    - 날씨 정보 조회

    - 모델 메타데이터 조회
\end{dhksmry}

등에 사용할 수 있다.

\textbf{POST}

서버에 데이터를 전달해서 처리하도록 요청할 때 주로 사용한다.

\begin{dhkcode}
response = requests.post(
    "https://api.example.com/predict",
    json={"income": 5000, "debt": 1000}
)
\end{dhkcode}

이 경우 서버에 데이터를 보내서 예측 결과를 요청하는 형태다.

여기서 한 가지 중요한 점은 \textbf{GET은 조회, POST는 생성이라고 외우면 안 된다는 것}이다.

POST는 새로운 데이터를 생성하는 데도 사용하지만, 데이터를 전달해서 계산이나 예측을 수행하도록 요청하는 데도 사용한다.

예를 들어 ML 추론 API가 대표적이다.

\begin{dhksmry}
    GET
    클라이언트 \ra 데이터 조회 요청 \ra 서버

    POST
    클라이언트 \ra 입력 데이터 전달 \ra 서버에서 처리
\end{dhksmry}

\subsubsection{Headers - 요청에 대한 부가 정보}

Header는 HTTP 요청이나 응답에 대한 메타데이터를 전달한다.

예를 들어:

\begin{dhkcode}
headers = {
    "Accept": "application/json"
}

response = requests.get(
    "https://api.example.com/data",
    headers=headers
)
\end{dhkcode}

여기서:

\begin{dhkzsh}
Accept: application/json
\end{dhkzsh}

은 클라이언트가 JSON 형식의 응답을 받을 수 있고 선호한다는 의미이다.

다만 서버가 반드시 JSON을 반환해야 한다는 뜻은 아니다.

다른 예시로:

\begin{dhkzsh}
headers = {
    "Authorization": "Bearer YOUR_API_TOKEN"
}
\end{dhkzsh}

같은 인증 정보도 Header에 넣을 수 있다.

이 부분은 Hugging Face나 외부 API를 사용할 때 자주 등장한다.

API 토큰은 실제 프로젝트에서 코드에 직접 적지 않고 환경변수 등으로 관리하는 것이 좋다.

\subsubsection{Body - 서버에 전달할 데이터}

이번에는 POST 요청을 보자.

\begin{dhkcode}
payload = {
    "income": 5000,
    "debt": 1000
}

response = requests.post(
    "https://api.example.com/predict",
    json=payload
)
\end{dhkcode}

여기서:

\begin{dhkzsh}
json=payload
\end{dhkzsh}

는 Python dictionary를 JSON 형식으로 직렬화하여

HTTP 요청의 Body에 넣도록 \zsh{requests}에 지시하는 것이다.

또한 일반적으로:

\begin{dhkzsh}
Content-Type: application/json
\end{dhkzsh}

이라는 Header도 자동으로 설정한다.

전체 구조를 보면:

\begin{dhkdir}[in=6]
HTTP Request
>Method: POST
>URL: /predict
>Headers
>>Content-Type: application/json
>Body
>>{"income": 5000, "debt": 1000}
\end{dhkdir}

여기서 JSON이 다시 등장한다.

우리는 이미:

\begin{dhkcode}
json.dumps(payload)
\end{dhkcode}

가 Python 객체를 JSON 문자열로 변환한다는 것을 배웠다.

\zsh{requests.post(...,json=payload)}는 이런 JSON 직렬화 과정을 내부적으로 처리해주는 것이다.

\subsubsection{HTTP REsponse는 무엇인가?}

클라이언트가 요청을 보내면 서버는 응답을 반환한다.

개념적으로:

\begin{dhkdir}[in=6]
HTTP Response
>Status Code
>>200 / 404 / 500 ...
>Headers
>>Content-Type 등
>Body
>>실제 응답 데이터
\end{dhkdir}

여기서 가장 먼저 살펴볼 것은 Status Code이다.

\textbf{Status Code} - 요청 처리 결과

HTTP 상태 코드는 서버가 요청을 어떻게 처리했는지 알려주는 숫자다.

\begin{dhktable}
|상태코드|의미|
|:---|:---|
|200|OK - 요청성공|
|201|Created - 리소스 생성 성공|
|400|Bad Request - 잘못된 요청|
|401|Unauthorized - 인증 필요 또는 인증 실패|
|403|Forbidden - 접근 거부|
|404|Not Found - 요청한 리소스를 찾을 수 없음|
|429|Too Many Requests - 요청 횟수 제한|
|500|Internal Server Error - 서버 내부 오류|
|503|Service Unavailable - 서버가 일시적으로 요청을 처리할 수 없음|
\end{dhktable}

개별 코드를 전부 암기할 필요는 없다.

대신 다음 분류를 이해하면 된다.

\begin{dhksmry}
    2xx \ra 성공

    3xx \ra 리다이렉션

    4xx \ra 클라이언트 측 요청 문제

    5xx \ra 서버 측 처리 문제
\end{dhksmry}

실제로는 인증, 권한, 요청 제한 등 여러 상황이 있으므로 4xx를 무조건 클라이언트 코드의 버그라고 생각하면 안 된다.

\subsubsection{\zsh{requests.Response}객체 이해하기}

이제 다음 코드를 보자.

\begin{dhkcode}
import requests

response = requests.get(
    "https://api.example.com/data",
    timeout=10
)
\end{dhkcode}

여기서 \zsh{response}는 단순한 문자열이 아니다.

\textbf{requests.Response} 객체이다.

이 객체는 서버에서 받은 응답 정보를 가지고 있다.

\begin{dhkcode}
response.status_code
response.headers
response.text
response.content
response.json()
\end{dhkcode}

각각 의미가 다르다.

\begin{dhktable}
|코드|반환 내용|
|:---|:---|
|response.status_code|HTTP 상태 코드(int)|
|response.headers|응답 Header 정보|
|response.text|응답 Body를 문자열로 디코딩한 결과(str)|
|response.content|응답 Body의 바이트 데이터(bytes)|
|response.json()|JSON Body를 Python 객체로 파싱한 결과|
\end{dhktable}

특히 다음 세 가지는 구분해야 한다.

\begin{dhkcode}
response.text
response.content
response.json()
\end{dhkcode}

앞에서 배운 File I/O와 연결하면:

\begin{dhkbranch}[in=8]
node main :: HTTP Response Body
branchline main -> rc :: response.content
arrow rc
node rc :: bytes
stop rc
branchline main -> rt :: response.text
arrow rt
node rt :: bytes
stop rt
handoff main -> json :: response.json()
arrow json
node json :: Python 객체 dict / list 등
stop json
\end{dhkbranch}

단, \zsh{response.json()}은 응답 본문이 유효한 JSON이어야 정상적으로 동작한다.

\textbf{JSON이 아닌 HTML이나 일반 텍스트를 반환했다면 파싱 오류가 발생할 수 있다}.

\subsubsection{\zsh{raise_for_status()}는 왜 필요할까?}

다음 코드를 보자.

\begin{dhkcode}
response = requests.get(url, timeout=10)
response.raise_for_status()
\end{dhkcode}

이 메서드는 HTTP 응답이 4xx 또는 5xx 오류 상태라면 예외를 발생시킨다.

예를 들어 서버가 404를 반환했다면:

\begin{dhkcode}
response.raise_for_status()
\end{dhkcode}

에서 \zsh{HTTPError}가 발생한다.

반면 정상적인 200 응답이라면 다음 코드로 진행할 수 있다.

\begin{dhkcode}
data = response.json()
\end{dhkcode}

여기서 중요한 차이가 있다.

\textbf{request.get()이 HTTP 응답을 받았다는 사실과 요청이 성공했다는 사실은 다르다}.

서버가 404 오류 페이지를 반환했더라도 HTTP 응답 자체는 수신했을 수 있다.

또한 네트워크 연결 자체가 실패하면 \zsh{Response}를 받기 전에 예외가 발생할 수도 있다.

그래서 기본적인 API 호출에서는:

\begin{dhkcode}
response = requests.get(url, timeout=10)
response.raise_for_status()
\end{dhkcode}

를 함께 사용하는 것이 좋다.

\subsubsection{공공데이터 API는 어떻게 호출할까?}

예를 들어 공공데이터 API가 다음과 같은 주소를 제공한다고 가정하자.

\begin{dhkzsh}
https://api.example.com/financial-products
\end{dhkzsh}

그리고 API 문서에 다음 파라미터가 정의되어 있다고 가정해보자.

\begin{dhksmry}
    page = 페이지 번호

    per_page = 페이지당 데이터 수
\end{dhksmry}

그러면 Python에서:

\begin{dhkcode}
import requests

url = "https://api.example.com/financial-products"
params = {
    "page": 1,
    "per_page": 100
}

response = requests.get(
    url,
    params=params,
    timeout=10
)

response.raise_for_status()
data = response.json()
\end{dhkcode}

처럼 요청할 수 있다.

여기서 새롭게 등장한 것은:

\begin{dhkcode}
params = params
\end{dhkcode}

이다.

이건 URL에 \textbf{Query Parameters}를 추가하는 기능이다.

실제 요청 URL은 개념적으로 다음과 같이 구성된다.

\begin{dhkzsh}
https://api.example.com/financial-products?page=1&per_page=100
\end{dhkzsh}

구조를 분해하면:

\begin{dhkdir}
https:/api.example.com/financial-products
> URL Path
\end{dhkdir}

\begin{dhkdir}
?page=1
> page = 1
\end{dhkdir}

\begin{dhkdir}
&per_page=100
> per_page = 100
\end{dhkdir}

즉 GET 요청에서는 필요한 조회 조건을 URL의 Query Parameters로 전달하는 경우가 많다.

하지만 모든 GET 요청이 반드시 Query Parameters를 사용해야 하는 것은 아니다.

\subsubsection{Query Parameters와 JSON Body의 차이}

둘을 비교해보자.

\textbf{GET 요청}

\begin{dhkcode}
requests.get(
    url,
    params={'page':1}
)

GET /financial-products?page=1
\end{dhkcode}

\textbf{POST 요청}

\begin{dhkcode}
requests.post(
    url,
    json={"income": 5000}
)

POST /predict

Body:
{"income": 5000}
\end{dhkcode}

여기에서 기억할 핵심은:

\begin{dhksmry}
    - params= 는 URL의 Query Parameters를 구성한다.

    - json= 은 JSON 형식의 요청 Body를 구성한다.
\end{dhksmry}

다만 Query Parameters는 POST에서도 사용할 수 있고,

POST Body가 반드시 JSON이어야 하는 것도 아니다

실제 사용 방식은 API 문서에서 정한다.

\subsubsection{공공데이터 API $\rightarrow$ Pandas 연결}

공공데이터 API가 다음과 같은 JSON을 반환한다고 가정하자.

\begin{dhkzsh}
{
    "data": [
        {
            "product_name": "상품A",
            "interest_rate": 3.5
        },
        {
            "product_name": "상품B",
            "interest_rate": 3.8
        }
    ]
}
\end{dhkzsh}

Python에서는:

\begin{dhkcode}
import requests
import pandas as pd

response = requests.get(url, timeout=10)
response.raise_for_status()

result = response.json()

records = result["data"]

df = pd.DataFrame(records)
\end{dhkcode}

처럼 처리할 수 있다.

전체 흐름을 보면:

\begin{dhkbranch}[in=12]
node main ::공공 데이터 API
arrow main :: HTTP GET
node main :: HTTP Response
arrow main :: response.json()
node main :: Python dict
arrow main :: result['data']
node main :: Python list[dict]
arrow main :: pd.DataFrame()
node main :: DataFrame
arrow main
\end{dhkbranch}

이렇게 Pandas와 연결된다.

실제 공공데이터 API에서는 인증키, 페이지네이션, 응답 형식 등이 서비스마다 다르기에

반드시 해당 API 문서를 확인해야 한다.

\subsubsection{다운로드와 파일 저장도 연결해보자}

마지막으로 API에서 받은 데이터를 파일로 저장하는 경우도 살펴보자.

\begin{dhkcode}
import requests
from pathlib import Path

url = "https://example.com/data/train.csv"

response = requests.get(url, timeout=10)
response.raise_for_status()

path = Path("data/train.csv")

path.parent.mkdir(
    parents=True,
    exist_ok=True
)

path.write_bytes(response.content)
\end{dhkcode}

여기서:

\begin{dhkcode}
response.content
\end{dhkcode}

는 \zsh{bytes} 객체이고,

\begin{dhkcode}
path.write_bytes(...)
\end{dhkcode}

는 해당 바이트 데이터를 파일에 기록하는 메서드다.

이 코드의 전체 흐름은:

\begin{dhksmry}[centered]
URL
\da
requests.get()
\da
Response object
\da
response.content
\da
bytes
\da
Path.write_bytes()
\da
Local file
\end{dhksmry}

가 된다.

이후:

\begin{dhkcode}
df = pd.read_csv(path)
\end{dhkcode}

로 읽으면 된다.

참고로 큰 파일은 \zsh{response.content}로 한 번에 메모리에 올리기보다

\zsh{stream=True}를 사용해서 나누어 저장하는 것이 적절할 수 있다.
